\chapter{اللقاء الثلاثون: الأيمان والإيلاء والفيء وبدايات أحكام العدة}

\section{مقدمة}
السلام عليكم ورحمة الله وبركاته. الحمد لله رب العالمين، وأشهد أن لا إله إلا الله وحده لا شريك له، وأشهد أن محمداً عبده ورسوله، صلى الله عليه وعلى آله وصحبه وسلم.

استأنف الشيخ هذا اللقاء بعد انقطاع يسير، فذكَّر بمقاصد مدارسة الكتب، وبأهمية الصبر على تراث الأئمة، ثم دخل في تفسير هذا الموضع من سورة البقرة، من قوله تعالى: \begin{ayah}لَا يُؤَاخِذُكُمُ اللَّهُ بِاللَّغْوِ فِي أَيْمَانِكُمْ\end{ayah} إلى أوائل أحكام الطلاق والعدة في قوله تعالى: \begin{ayah}وَالْمُطَلَّقَاتُ يَتَرَبَّصْنَ بِأَنْفُسِهِنَّ ثَلَاثَةَ قُرُوءٍ\end{ayah}.

وكان هذا المجلس غنياً جداً من جهة المنهج؛ إذ ظهر فيه صنيع \rawi{الطبري} في توجيه الأقوال، وشرح مأخذ كل قول، والتمييز بين مواضع الاتفاق ومواضع النزاع، ثم ختم ذلك بترجيح محكم يجمع بين لسان العرب والآثار والسياق.

\separator

\section{تأويل (لا يؤاخذكم الله باللغو في أيمانكم)}
\isnad{قال أبو جعفر:} بيَّن الله في هذه الآية أنه لا يؤاخذ عباده باللغو من أيمانهم، ثم ذكر \rawi{الطبري} أقوال السلف في معنى هذا اللغو، ووازن بينها بميزان دقيق:
\begin{itemize}
    \item فقيل: هو ما يجري على اللسان من غير قصد عقد اليمين.
    \item وقيل: هو يمين المخطئ الذي يحلف يظن الأمر كما قال ثم يتبين له خلافه.
    \item وقيل: يدخل فيه بعض ما لا يقصد به الالتزام ابتداء.
\end{itemize}

ثم أبطل ما لا يستقيم مع دلالة الآية، وقرر أن اللغو هو ما لا يقصد صاحبه به عقد اليمين على الوجه الذي يترتب عليه حكم المؤاخذة. ولهذا كان من أفسد الأقوال عنده أن يوصف شيء بأنه لغو، ثم توجب فيه الكفارة؛ لأن اللغو الذي نفى الله المؤاخذة به لا يجتمع مع إلزام صاحبه حكم المؤاخذة.

\begin{fawaid}[تحرير معنى اللغو يرفع إشكالاً كبيراً في باب الأيمان]
من أجمل ما في هذا الموضع أن \rawi{الطبري} لا يكتفي بسرد الأقوال، بل يختبر كل قول على ضوء الآية نفسها: فما نفى الله المؤاخذة به لا يصح أن يثبت له في الوقت نفسه حكمها، وإلا صار التفسير متناقضاً.
\end{fawaid}

\section{تأويل (ولكن يؤاخذكم بما كسبت قلوبكم)}
بعد تقرير معنى اللغو انتقل \rawi{الطبري} إلى القدر المقابل له، وهو ما تكسبه القلوب من الأيمان، فذكر أن هذا الكسب هو القصد والعزم عن علم ومعرفة، لا مجرد لفظ اللسان.

ثم لخص المعنى الراجح في صورتين:
\begin{itemize}
    \item صورة يحلف فيها المرء على أمر ماضٍ أو واقع، وهو يعلم أنه كاذب، قاصداً الكذب والباطل.
    \item وصورة يعقد فيها اليمين على أمر مستقبل، ثم يحنث بعد ذلك فيما حلف عليه.
\end{itemize}

أما الأولى فهي من اليمين الآثمة التي يتوعد صاحبها في الآخرة، ولا كفارة فيها في الدنيا عند تقرير \rawi{الطبري} هنا؛ لأنها ليست من جنس اليمين التي يدخلها الحنث بعد العقد. وأما الثانية فهي اليمين المعقودة على المستقبل، وهذه إذا حنث صاحبها لزمته الكفارة.

وهنا تظهر فائدة دقيقة: أن المؤاخذة ليست نوعاً واحداً، بل قد تكون مؤاخذة أخروية بالإثم والوعيد، وقد تكون مؤاخذة دنيوية بإيجاب الكفارة، بحسب نوع اليمين ومحلها.

\begin{fawaid}[التمييز بين اليمين على الماضي واليمين على المستقبل أصل مهم]
ليس كل يمين تعمدها القلب يجري عليها الحكم نفسه؛ فاليمين الكاذبة على الماضي بابها الإثم والوعيد، واليمين المعقودة على المستقبل بابها الحنث والكفارة. وإذا لم يحرر هذا الفرق اختلطت أبواب الأيمان على القارئ.
\end{fawaid}

\section{تفسير خاتمة الآية: (والله غفور حليم)}
وقف الشيخ عند تفسير \rawi{الطبري} لخاتمة الآية، فذكر أنه فسر اسمي الله هنا بما يناسب المقام:
\begin{itemize}
    \item فهو \textit{غفور} لعباده فيما جرى من لغو أيمانهم.
    \item وهو \textit{حليم} لا يعاجل أهل المعصية بالعقوبة، مع قدرته عليهم.
\end{itemize}

وهذا من دقائق التفسير؛ فإن أسماء الله في خواتيم الآيات ليست زوائد لفظية، بل هي مفاتيح لفهم السياق، وتكميل المعنى، وربط الحكم بصفة الرب سبحانه.

\begin{fawaid}[خواتيم الآيات من تمام التفسير]
من نقص التفسير أن تذكر الحكم وتغفل عن الخاتمة الإلهية، لأن اقتران الحكم باسم الله أو صفته يبين وجه التشريع، ويزيد القلب فهماً وتعظيماً.
\end{fawaid}

\separator

\section{تأويل (للذين يؤلون من نسائهم تربص أربعة أشهر)}
\isnad{قال أبو جعفر:} الإيلاء في أصله الحلف، ثم صار اسماً مخصوصاً في هذا الباب للحلف الذي يمتنع به الزوج من قربان زوجته.

وعرض \rawi{الطبري} أقوال أهل العلم في حد الإيلاء:
\begin{itemize}
    \item فقيل: هو كل حلف على ترك جماع الزوجة.
    \item وقيل: إنما يكون إذا كان ذلك في حال غضب أو إضرار.
    \item وقيل: العبرة بوقوع الحلف على ترك الوطء مدةً يتحقق بها الضرر، لا بمجرد وصف الغضب.
\end{itemize}

ثم رجح أن الإيلاء المراد في الآية هو الحلف المتعلق بترك الجماع، إذا كان على وجه يمنع حق الزوجة وتترتب عليه مدة التربص المذكورة، سواء كان ذلك في حال غضب أو في غيره. فالعبرة عنده بحقيقة الامتناع الذي انعقد عليه اليمين، لا بالقيد الذي لم يدل عليه نص الآية.

وفي هذا الموضع ظهر أيضاً منهج \rawi{الطبري} في شرح الأقوال لا مجرد حكايتها؛ فهو يبين كيف فهم كل فريق الآية، ولماذا وسع بعضهم المعنى أو ضيقه، ثم يرد الجميع إلى القدر الذي يدل عليه اللفظ والسياق.

ونبّه الشيخ إلى أن من أهم ما يضبط هذا الباب معرفة محل الاتفاق ومحل النزاع؛ فالجميع يتكلمون عن يمين ترتب ضرراً على الزوجة، لكن الخلاف وقع في بعض القيود: هل كل امتناع يدخل في الإيلاء، أم لا بد أن يكون الحلف على ترك الجماع نفسه؟ وهذا التحرير هو الذي تتفرع عنه بقية المسائل.

\begin{fawaid}[شرح الأقوال قبل الحكم عليها من خصائص التفسير المحكم]
من الإنصاف العلمي أن يفهم القارئ مراد كل قول قبل أن يحكم له أو عليه. وهذه طريقة ظاهرة عند \rawi{الطبري}: يفسر الأقوال أولاً، ثم يبين ما يدخل في الآية وما يخرج عنها.
\end{fawaid}

\section{تأويل (فإن فاؤوا فإن الله غفور رحيم)}
لما ذكر الله التربص في الإيلاء عقبه بذكر الفيء، أي: الرجوع عما حلف عليه الزوج من الامتناع. وهنا نقل \rawi{الطبري} خلاف السلف في حقيقة الفيء:
\begin{itemize}
    \item فقيل: الفيء هو الرجوع إلى الجماع نفسه.
    \item وقيل: قد يقوم مقامه القول أو العزم إذا كان الزوج عاجزاً.
\end{itemize}

والذي يترجح من صنيعه في هذا الموضع أن الأصل في الفيء هو الرجوع إلى الوطء، لأنه المقابل للامتناع الذي وقع عليه الإيلاء ابتداء، وبه يعود الحق إلى موضعه الأصلي.

ومن دقائق تحريره هنا أن الخلاف في معنى الفيء تابع للخلاف في معنى الإيلاء نفسه؛ فمن جعل الإيلاء مخصوصاً بالحلف على ترك الجماع جعل الفيء رجوعاً إلى الجماع، ومن وسع الإيلاء إلى أنواع أخرى من الامتناع أو الإضرار وسع معنى الفيء تبعاً لذلك. ثم رجح \rawi{الطبري} أن العاجز عن الجماع يكفيه عزم القلب وإظهار الرجوع حتى يقدر، أما القادر فلا يتحقق فيؤه إلا بالفعل.

ولذلك كانت مناسبة الختم بقوله: \begin{ayah}فَإِنَّ اللَّهَ غَفُورٌ رَحِيمٌ\end{ayah} ظاهرة؛ لأن من رجع إلى ما أذن الله به، وترك الإضرار، دخل في باب المغفرة والرحمة.

\section{تأويل (وإن عزموا الطلاق فإن الله سميع عليم)}
بعد بيان الفيء انتقل \rawi{الطبري} إلى ما يقع إذا لم يرجع الزوج، وهنا فصل الخلاف تفصيلاً نافعاً:
\begin{itemize}
    \item ففريق رأى أن مضي الأربعة الأشهر نفسه يوقع الطلاق.
    \item ثم اختلف هؤلاء: هل تقع به طلقة بائنة أم رجعية؟
    \item وفريق آخر رأى أن مضي المدة لا يوقع الطلاق بنفسه، بل يوجب إيقاف الزوج وتخييره بين الفيء والطلاق.
\end{itemize}

والذي رجحه \rawi{الطبري} أن الزوج لا تبين منه امرأته بمجرد انقضاء الأربعة الأشهر، بل يوقف بعد مضيها، فإن فاء أو طلق عمل بما اختاره. ومن أجمل ما استدل به هنا نظره إلى خاتمة الآية نفسها: فإن قوله تعالى \begin{ayah}فَإِنَّ اللَّهَ سَمِيعٌ عَلِيمٌ\end{ayah} أنسب بما يسمع من طلاق ينطق به الزوج، لا بمجرد انقضاء مدة لا توصف بأنها مسموعة.

\begin{fawaid}[خاتمة الآية قد ترجح بين احتمالين فقهيين]
ليس استحضار أسماء الله وصفاته في أواخر الآيات أمراً تكميلياً فقط، بل قد يكون من أدلة الترجيح نفسها. وقد استعمل \rawi{الطبري} هنا ختم الآية بـ \textit{سميع عليم} ليرجح أن الطلاق لا يقع بمجرد مضي الزمن، بل عند العزم الذي يظهر أثره في القول والفعل.
\end{fawaid}

\separator

\section{تأويل (والمطلقات يتربصن بأنفسهن ثلاثة قروء)}
دخل الشيخ بعد ذلك في أول أحكام العدة، فوقف مع لفظ \textit{القروء}، وهو من الألفاظ التي اشتهر فيها الخلاف بين السلف والفقهاء:
\begin{itemize}
    \item فقيل: المراد بها الحيض.
    \item وقيل: المراد بها الأطهار.
\end{itemize}

وبيَّن \rawi{الطبري} سبب الاشتباه هنا من جهة اللسان؛ فإن القرء في كلام العرب يطلق على الإقبال والإدبار، ولذلك استعمل فيما يتصل بالحيض والطهر جميعاً. ومن هنا نشأ الخلاف، لا من مجرد التشهي.

ثم بنى على ذلك ترجيحه، فذكر أن الطلاق المشروع إنما يكون في طهر لم يجامعها فيه، وأن العدة المأمور بها هي ثلاثة قروء تأتي بعد هذا الطلاق على الوجه الذي دل عليه ظاهر التنزيل. ومن مجموع تقريره في هذا الموضع يظهر أن مذهبه يميل إلى أن القرء هنا هو \textit{الطهر}، وأن انقضاء العدة يكون بانقضاء الطهر الثالث وابتداء الحيضة التي تليه.

\begin{fawaid}[معرفة سبب الخلاف تعين على حسن الترجيح]
لا يكفي أن يقال: اختلفوا في القرء هل هو الحيض أو الطهر، بل من المهم أن يبين سبب هذا الاختلاف في لسان العرب واستعماله، لأن فهم منشأ النزاع أدعى إلى العدل في النظر وأقرب إلى إصابة الترجيح.
\end{fawaid}

\section{أثر هذا الترجيح في باب العدة}
نبَّه الشيخ إلى أن هذا الترجيح ليس مسألة نظرية محضة، بل ينبني عليه تحديد انتهاء العدة عملياً. فإذا كان القرء هو الطهر، وكانت المرأة قد طلقت في طهر لم يجامعها فيه، فإن عدتها تستكمل بالأطهار الثلاثة، فإذا انقضى الطهر الثالث ودخلت في الحيض الذي بعده، خرجت من العدة.

وهذا الموضع من أحسن أمثلة الجمع بين دلالة اللفظ، والحكم الفقهي، وصورة المسألة الواقعة، وهو ما يجعل قراءة \rawi{الطبري} نافعة للمفسر والفقيه معاً.

\separator

وقف الشيخ في ختام اللقاء عند هذا الموضع من الآية، وأخر الكلام على قوله تعالى: \begin{ayah}وَلَا يَحِلُّ لَهُنَّ أَنْ يَكْتُمْنَ مَا خَلَقَ اللَّهُ فِي أَرْحَامِهِنَّ\end{ayah} إلى المجلس التالي.

وخلاصة هذا اللقاء أن \rawi{الطبري} جمع فيه بين دقة الفقه، وتحقيق اللسان، وحسن عرض الأقوال، وربط الأحكام بخواتيم الآيات، فخرج القارئ منه بأصول نافعة في فهم الأيمان، والإيلاء، والفيء، وبدايات أحكام العدة.
